% Template: Đồ thị hai hàm bậc nhất giao nhau
% y1 = 2x - 1, y2 = -x + 3, giao tại G(4/3; 5/3)
% --------------------------------------------------------
\documentclass[border=5pt]{standalone}
\usepackage[utf8]{inputenc}
\usepackage[T5]{fontenc}
\usepackage[vietnamese]{babel}
\usepackage{tikz}
\usepackage{helvet}
\renewcommand{\familydefault}{\sfdefault}
% ── Hệ màu tập trung (chỉnh tại đây để đổi theme) ──────────
\usepackage{xcolor}
\definecolor{mauchinh} {HTML}{1565C0}
\definecolor{mauphu}   {HTML}{43A047}
\definecolor{maunoibat}{HTML}{E53935}
\definecolor{mauvien}  {HTML}{424242}
\usetikzlibrary{calc, arrows.meta}

\begin{document}
\begin{tikzpicture}[
    % Khai báo hai hàm số bằng declare function
    declare function={
        % Hàm bậc nhất thứ nhất: y = 2x - 1
        f(\x) = 2*\x - 1;
        % Hàm bậc nhất thứ hai: y = -x + 3
        g(\x) = -\x + 3;
    },
    scale=0.9,
    >=Stealth,
    line join=round,
    line cap=round,
    every node/.style={font=\sffamily\small, mauvien},
]
    % Giới hạn vẽ
    \pgfmathsetmacro{\xmin}{-1}
    \pgfmathsetmacro{\xmax}{3.5}
    \pgfmathsetmacro{\ymin}{-2}
    \pgfmathsetmacro{\ymax}{4}
    % x giao = 4/3 (từ 2x-1 = -x+3)
    \pgfmathsetmacro{\xgiao}{4/3}
    % y giao = 5/3
    \pgfmathsetmacro{\ygiao}{5/3}

    % Trục Ox
    \draw[mauvien, ->] (\xmin-0.3, 0) -- (\xmax+0.3, 0) node[right] {$x$};
    % Trục Oy
    \draw[mauvien, ->] (0, \ymin-0.3) -- (0, \ymax+0.3) node[above] {$y$};
    % Gốc tọa độ O
    \path (0,0) node[below left] {$O$};

    % Vạch và nhãn trục x
    \foreach \x in {-1,1,2,3} {
        \draw[mauvien] (\x, 2pt) -- (\x, -2pt);
        \path (\x,0) node[below] {$\x$};
    }

    % Vạch và nhãn trục y
    \foreach \y in {-1,1,2,3} {
        \draw[mauvien] (2pt,\y) -- (-2pt,\y);
        \path (0,\y) node[left] {$\y$};
    }

    % Đường thẳng y = f(x) = 2x - 1
    \draw[mauchinh, thick]
        plot[domain=\xmin:\xmax, variable=\x] (\x, {f(\x)});
    % Nhãn d1
    \path (\xmax, {f(\xmax)}) node[mauchinh, right] {$y = 2x-1$};

    % Đường thẳng y = g(x) = -x + 3
    \draw[maunoibat, thick]
        plot[domain=\xmin:\xmax, variable=\x] (\x, {g(\x)});
    % Nhãn d2
    \path (\xmin, {g(\xmin)}) node[maunoibat, left] {$y = {-x}+3$};

    % Giao điểm G(4/3; 5/3)
    \path (\xgiao, \ygiao) coordinate (G);
    \fill[mauphu] (G) circle (2.5pt);
    % Nhãn giao điểm
    \path (G) node[above right] {$G\!\left(\frac{4}{3};\frac{5}{3}\right)$};

    % Đường chiếu giao điểm xuống Ox
    \draw[mauvien!40, dotted, thin] (\xgiao, 0) -- (G);
    % Đường chiếu giao điểm sang Oy
    \draw[mauvien!40, dotted, thin] (0, \ygiao) -- (G);
\end{tikzpicture}
\end{document}
